\chapter{Mercadotecnia digital: la campaña}
\label{cap:mkt}
\textit{Criterio 7 de la rúbrica (8\,\%): estrategia integral sustentada en datos; plan mercadológico, branding, buyer
persona y medios digitales.}

\section{El marco: segmentar, elegir y posicionar}
La campaña sigue el esquema clásico de mercadotecnia \textbf{STP}:
\begin{enumerate}
\item \textbf{Segmentar}: dividir a los visitantes posibles en grupos con datos (de dónde vienen, cuándo viajan, cómo se
  informan).
\item \textbf{Elegir} (\emph{targeting}): a qué grupos hablarles.
\item \textbf{Posicionar}: qué lugar quiere ocupar la marca en su mente frente a las alternativas.
\end{enumerate}
La diferencia con una campaña común es que cada paso sale de una cifra, y cada rasgo dice si es dato, cálculo o supuesto.

\section{Segmentar: lo que dicen los datos}
\begin{itemize}
\item \textbf{Quién visita el sur} (INAH 2025): Oxtankah, 94.8\,\% nacionales; la Ruta, 63.3\,\% nacionales y 36.7\,\%
  extranjeros.
\item \textbf{Quién llega al estado en avión} (Unidad de Política Migratoria, 2025): 9,408,423 extranjeros a Cancún;
  Estados Unidos 56.3\,\% y Canadá 17.1\,\%; 53.4\,\% mujeres; más llegadas en marzo, menos en septiembre.
\item \textbf{El hallazgo que define todo}: al aeropuerto de Chetumal llegaron \textbf{250} extranjeros en 2025. El
  extranjero no llega al sur por avión: entra por Cancún y baja por carretera, en Tren Maya, desde Belice o en crucero.
\item \textbf{Cómo se informan} (ENDUTIH 2025, INEGI): 92.1\,\% de la población de Quintana Roo usa internet (tercer lugar
  nacional); de quienes tienen teléfono inteligente, 90.6\,\% usa mensajería y 80.4\,\% redes sociales.
\item \textbf{Qué valoran y qué los aleja} (85,987 reseñas, capítulo~\ref{cap:mineria}): el ruido (2.76$\times$), la
  suciedad (1.89$\times$), el precio (1.66$\times$) y las multitudes (1.52$\times$) hunden una reseña; la calma
  (0.41$\times$) y la cultura (0.48$\times$) la protegen.
\end{itemize}

\section{Elegir: dos viajeras ideales (buyer persona)}
\begin{decision}
Dos personas: ``La que vuelve al sur'' (nacional) y ``La que baja del norte'' (Estados Unidos o Canadá, ya en Cancún o la
Riviera). Descartadas: solo la nacional (deja fuera al extranjero que ya está a 334 km) y solo la extranjera (la Bahía es
95\,\% nacional).
\end{decision}

\small
\begin{longtable}{P{0.19\linewidth}P{0.17\linewidth}P{0.42\linewidth}P{0.1\linewidth}}
\rowcolor{guinda}\h{Persona} & \h{Rasgo} & \h{Valor y fuente} & \h{Tipo} \\ \endhead
\textbf{La que vuelve al sur} & Origen & Nacional: 94.8\,\% de Oxtankah y 63.3\,\% de la Ruta (INAH 2025) & dato \\
 & Cuándo viaja & De diciembre a abril; la campaña la invita en los meses con espacio (forma del año) & derivado \\
 & Qué valora & Calma (0.41$\times$) y cultura (0.48$\times$) en las reseñas & derivado \\
 & Qué la aleja & Ruido, suciedad, precio y multitudes & derivado \\
 & Cómo se informa & Mensajería 90.6\,\%, redes 80.4\,\% (ENDUTIH 2025) & dato \\
 & Estado de origen, edad, ingreso & No hay dato oficial por destino & hueco \\
\textbf{La que baja del norte} & Origen & De 9,408,423 extranjeros en Cancún: EE.\,UU. 56.3\,\%, Canadá 17.1\,\% & dato \\
 & Sexo & 53.4\,\% mujeres & dato \\
 & Por qué no llega sola & 250 extranjeros en el aeropuerto de Chetumal & dato \\
 & Cuándo está & Más en marzo, menos en septiembre; Cancún en temporada alta de noviembre a abril & dato \\
 & Cómo se informa & Busca en Google y usa redes desde el destino & supuesto \\
 & Idioma & Inglés y español & derivado \\
\end{longtable}
\normalsize

\section{Posicionar: la marca}
\begin{itemize}
\item \textbf{Nombre}: ``El sur tiene espacio'' (``The south has room''). Dice el dato (capacidad ociosa) y la promesa (sin
  multitudes) a la vez.
\item \textbf{Lema}: ``Pirámides, bahía y sabor del sur, con espacio para disfrutarlos''.
\item \textbf{Personalidad} (decisión de Brandon: combinar los tres tonos, ``porque es un orgullo ser mexicano''): cercana y
  tranquila; orgullosa de lo maya y lo mexicano; con espíritu de aventura.
\item \textbf{Propuesta de valor}: el Caribe mexicano que todavía tiene espacio: zonas mayas con 15.5 veces menos visitantes
  que Tulum (INAH 2025), una bahía tranquila y la comida del sur.
\item \textbf{Posicionamiento}: para quien quiere el Caribe mexicano sin multitudes, el sur de Quintana Roo es la opción
  cultural y tranquila; a diferencia de Tulum o Cancún, aquí hay espacio, y la campaña lo cuida.
\item \textbf{Elementos visuales}: el sistema ``Sur mexicano'' (rosa, amarillo, turquesa y añil en bloques; letras Bricolage
  y Figtree; greca maya; fotos reales con coordenada comprobada).
\end{itemize}

\section{Los mensajes: cada promesa con su dato}
Reglas: cada anuncio lleva el dato que lo respalda; \textbf{sin precios ni horas de viaje} (no hay dato oficial); el
lenguaje sale de las reseñas de 5 estrellas; el largo se revisa contra los límites de cada plataforma (Google: título
$\le30$ caracteres y descripción $\le90$; Meta: título $\le40$ y texto $\le125$).

\begin{center}
\small
\begin{tabularx}{\linewidth}{llY}
\encabezado \h{Persona} & \h{Canal} & \h{Mensaje y respaldo} \\
Vuelve & Google (es) & ``Pirámides mayas con espacio''. \emph{Respaldo}: Tulum 1,031,443 visitantes; Ruta 66,628 (15.5 veces menos). \\
Vuelve & Facebook (es) & ``La bahía que te esperaba''. \emph{Respaldo}: Oxtankah, 11,017 visitantes en 2025, 94.8\,\% mexicanos. \\
Vuelve & WhatsApp & ``¿Y si este puente nos vamos al sur?''. \emph{Respaldo}: 90.6\,\% usa mensajería; sin costo. \\
Baja & Google (en) & ``Mayan pyramids, with room''. \emph{Respaldo}: el mismo 15.5; Cancún en temporada alta. \\
Baja & Facebook (en) & ``Your Caribbean trip, one stop south''. \emph{Respaldo}: Tren Maya con estación en Chetumal; solo se muestra si la Torre marca el norte lleno. \\
\bottomrule
\end{tabularx}
\normalsize
\end{center}

\section{Los medios y el calendario}
\begin{center}
\begin{tabularx}{\linewidth}{lrY}
\encabezado \h{Medio} & \h{Peso} & \h{Justificación con datos} \\
Facebook / Instagram & 70\,\% & 6.61 visitantes por cada \$1,000; 80.4\,\% usa redes. Tope de 70\,\% para no depender de una plataforma \\
Google búsqueda & 30\,\% & Llega a quien ya busca viajar; conversión medida de 5.75\,\% en viajes \\
WhatsApp (orgánico) & sin costo & 90.6\,\% usa mensajería; el botón de compartir viaja entre familia y amigos \\
La página & destino & Planeador por lugar y mes en 9 idiomas, sin llamadas a otros sitios \\
\bottomrule
\end{tabularx}
\end{center}
\textbf{Calendario}: los 9 meses del plan (octubre de 2026 a junio de 2027) con el dinero de la Fase 6. ``La que baja del
norte'' se activa solo en los meses en que Cancún está en temporada alta y el sur tiene espacio; además, la Torre la
enciende por semana cuando el norte se satura.

\section{Cómo se mide: los indicadores}
\textbf{El embudo con números} (con las tasas de referencia de viajes, y con el reparto de la Fase 6):
\begin{amano}[ · impresiones, clics y visitantes esperados]
Facebook: $131{,}250$ pesos $\div\ 8.70$ pesos por clic $=15{,}084$ clics. Con una tasa de clics de 2.76\,\%, eso supone
$15{,}084/0.0276=546{,}535$ impresiones, y con conversión de 5.75\,\%, 867 visitantes.

Google: $56{,}250\div36.17=1{,}555$ clics; con tasa de clics de 8.73\,\%, 17,814 impresiones; 89 visitantes.

Son estimaciones con promedios de EE.\,UU.: las metas reales se fijan con el primer mes de campaña.
\end{amano}

\begin{center}
\small
\begin{tabularx}{\linewidth}{lYY}
\encabezado \h{Indicador} & \h{Fórmula} & \h{Meta} \\
Alcance & Personas únicas que vieron un anuncio & Se fija con el primer mes real \\
Interacción (CTR) & clics $\div$ impresiones & Google $\ge8.73\,\%$; Facebook $\ge2.76\,\%$ \\
Conversión & planes de viaje $\div$ clics & $\ge5.75\,\%$ \\
Costo por visitante & pesos gastados $\div$ conversiones & $\le\$196$ \\
Afluencia & visitantes reales $-$ escenario probable & +957 en 9 meses \\
Presupuesto & pesos gastados $\div$ pesos planeados & 100\,\% (lo pausado se gasta después) \\
Económico & ocupación hotelera de Chetumal & \textbf{Hueco}: no se publica desde 2025 \\
Social & visitantes por cada mil habitantes (Radar) & Nunca ``saturado'' \\
Ambiental & semanas con anuncio en temporada alta o con mal clima & 0 \\
Capacidad & (esperados + campaña) $\div$ capacidad probada & $\le40\,\%$ en meses con anuncio \\
\bottomrule
\end{tabularx}
\normalsize
\end{center}

\section{La regla anti-colapso}
El incidente de mercadotecnia pregunta qué hacer si la afluencia rebasa la capacidad. La campaña tiene la respuesta
programada, no solo escrita:
\begin{enumerate}
\item Nunca anuncia un lugar en su temporada alta (restricción del presupuesto).
\item No permite que los visitantes esperados más los de la campaña pasen de la capacidad probada.
\item La Torre pausa el anuncio si hay tormenta, clima raro o llegó más gente de la esperada.
\item A quien iba al norte lleno le ofrece el lugar del sur con más espacio, nunca otro lugar lleno.
\end{enumerate}

\begin{codigoen}\codigo{campana/marca.py} (\codigo{personas}, \codigo{marca}, \codigo{mensajes}, \codigo{revisar_largos},
\codigo{medios}, \codigo{calendario}, \codigo{kpis}) y \codigo{campana/texto.py}. Notebook \codigo{07_campana}.\end{codigoen}
